\documentclass[10pt,letterpaper]{article}
\usepackage{cornell-notes}

\title{Clause 03.15: Model Kind}
\author{}
\date{}
\setDocTitle{Clause 03.15: Model Kind}
\setDocSubtitle{ISO/IEC/IEEE 42010:2022}
\setDocOwner{ISO 42010 Cornell Notes}
\setDocDate{\today}
\setCornellCollection{ISO/IEC/IEEE 42010:2022 Cornell Notes}
\setCornellUnitType{Clause}
\setCornellUnitNumber{03.15}
\setCornellUnitTitle{Model Kind}
\hypersetup{pdftitle={Clause 03.15: Model Kind},pdfsubject={Cornell notes},pdfauthor={}}

\begin{document}
\maketitle
\begin{CornellOverviewBox}{Overview}
{\Large\bfseries\textcolor{CNavy}{Section 3.15: model kind}}\\[3pt]
{\small\textbf{Classification:} Definition}\\
{\small\textbf{Learning objective:} Explain the precise meaning of model kind and distinguish it from closely related architecture-description concepts.}
\end{CornellOverviewBox}
\vspace{3pt}

\begin{CornellNotesTable}
\CornellNoteRow{Term}{\textbf{model kind}}
\CornellNoteRow{Definition}{A category of model distinguished by defining characteristics and modelling conventions.}
\CornellNoteRow{Distinction}{A model kind defines a category and modelling conventions; a view component is an actual separable portion of a view governed by a model kind or legend.}
\CornellNoteRow{Example / application}{Functional models, activity models, structural models, use case models, geopolitical models, analytic models and economic models.}
\CornellNoteRow{Active recall}{\begin{itemize}[leftmargin=*,nosep,topsep=1pt]
\item How would you explain model kind in your own words?
\item Which neighboring concept is easiest to confuse with model kind, and how is it different?
\end{itemize}}
\end{CornellNotesTable}


\begin{CornellSummaryBox}{Summary}
A category of model distinguished by defining characteristics and modelling conventions. A model kind defines a category and modelling conventions; a view component is an actual separable portion of a view governed by a model kind or legend.
\end{CornellSummaryBox}

\end{document}
